\icode{E-EXT-FI-L1-260625-001}
在用油膜法估测油酸分子的大小的实验中，具体操作如下：
  \begin{center}
  \begin{minipage}[c]{.9\textwidth}
    \ding{172}取油酸\,\qty{1.0}{mL}注入\,\qty{250}{mL}的容量瓶内，然后向瓶中加入酒精，直到液面达到\,\qty{250}{mL}的刻度为止，摇动瓶使油酸在酒精中充分溶解，形成油酸酒精溶液；用滴管吸取制得的溶液逐滴滴入量筒，记录滴入的滴数直到量筒达到\,\qty{1.0}{mL}为止，恰好共滴了100滴；


  \ding{173}将玻璃板放在浅盘上，然后将油膜的形状用彩笔描绘在玻璃板上。

  \ding{174}用注射器将配好的油酸酒精溶液滴一滴在水面上，待油酸膜形状稳定。

  \ding{175}将画有油膜形状的玻璃板放在边长为\,\qty{1}{cm}的方格纸上，算出完整的方格有67个，大于半格的有14个，小于半格的有19个。

  \ding{176}往边长约为\,\qty{40}{cm}的浅盘里倒入约\,\qty{2}{cm}深的水，待水面稳定后将适量的痱子粉均匀地撒在水面上。
  \end{minipage}
  \end{center}

  （1）上述步骤中，正确的顺序是\underline{\hbox to 7em{}}

  （2）利用上述具体操作中的有关数据可知一滴油酸酒精溶液含纯油酸为\underline{\hbox to 5em{}}\,\qty{}{m^3}，油膜面积为\underline{\hbox to 5em{}}\,\qty{}{m^2}，求得的油膜分子直径为\underline{\hbox to 5em{}}\,\qty{}{m}。（结果全部取2位有效数字）